Utilities reviewing large grid-project queues need portfolio-aware optimization together with a record of the interventions that shape each recommendation. TRACE-MOEA combines preference-guided ranking, deterministic budget repair, and a quarantined decision trace that records repair drops and preference-elite injections without entering the evaluation metric. On a reproducible five-objective benchmark of 120 candidate projects, the method reaches pooled mean hypervolume 0.17425 across seven review scenarios and 30 seeds per stochastic method, 0.89\% above NSGA-II (0.17270). Among 28 comparisons with four stochastic baselines, TRACE-MOEA has 27 positive mean differences and 24 Holm-significant wins, with no significant loss; 21 gaps against three deterministic ranking rules are descriptive and all favor TRACE-MOEA. A separate 270-run budget scan yields higher hypervolume than R-NSGA-II and NSGA-II at all three tested budgets. Component attribution is narrower: removing preference adaptation changes pooled hypervolume by 0.17\%, and the direct contrast is unresolved after cross-scenario correction. The decision trace covers 98.6\% of projects in returned fronts, but its effect on human review remains untested. Public reliability and project-outcome checks provide descriptive external consistency rather than validated review performance.
